% ===================================================================
%  GAMBAR No. 16 — Toko Gedhe. --- Pasar. Dolanan.
%
%  Traduction, faite en 2026, en javanais d'aujourd'hui, sur l'ido de
%  texto/io. Regles de la colonne : voir 10-gambar-01.tex. Une seule
%  numerotation. Dernier tableau du livret.
%
%  ECART DECLARE, ET IL FAUT LE SAVOIR AVANT D'ECRIRE. Le bloc
%  t16-15-1 est l'un des trois que renvoji.py laisse passer sans rien
%  dire : le fac-simile ido saute la lettre « (d) » des hemispheres,
%  que la planche grave pourtant. Un ecart declare rend le bloc
%  AVEUGLE — le controle n'y regarde plus rien du tout. L'ordre
%  85 a b c d e f 86 87 88 89 a donc ete pose a la main et relu a la
%  main, comme au tableau 8 pour le gobelet (13).
%
%  LE wayang EST REFUSE, ET C'EST LE REFUS QUE TOUTE LA COLONNE
%  PREPARAIT. Le theatre de marionnettes (41) est un petit theatre
%  europeen a rideau rouge, decors, loges et fauteuils d'orchestre. Le
%  javanais a « wayang », qui est le mot le plus javanais du monde et
%  qui nomme un theatre de marionnettes ; lui donner ce nom aurait
%  modernise LA CHOSE. La colonne ecrit « panggung boneka ». C'est le
%  septieme refus, apres gobak sodor (tableau 2), joged (3),
%  gendhongan (6), gamelan (11) et les jeux du cafe (13) — et le seul
%  qui coute vraiment quelque chose, parce que wayang est le seul des
%  sept dont la chose ressemble a s'y meprendre a la chose gravee.
%  C'est justement pourquoi il fallait le refuser.
%
%  MAIS LE VIOLON (76) EST ACCEPTE, ET SOUS LE MEME NOM PORTUGAIS
%  QU'EN INDONESIEN : biola, de viola. La colonne ourdoue avait refuse
%  l'emprunt et garde سارنگی ; les deux colonnes de Java le prennent.
%  Le mot suit l'objet quand l'objet est venu, et le refuse quand la
%  langue en avait deja un : le javanais avait le rebab et le gamelan,
%  mais pas le violon d'orchestre, qui est arrive avec ses navires. Un
%  instrument n'est pas un autre instrument parce qu'on en joue avec
%  un archet.
%
%  LES SERDADU TIMAH (29) SONT PORTUGAIS COMME EN INDONESIEN, ET LES
%  VRAIS SOLDATS DU § 4 SONT SANSKRITS — prajurit. Le jouet garde le
%  vieux mot, la chose garde le sien : la meme separation qu'au
%  tableau 6 entre le pawon et le kompor, mais tiree cette fois entre
%  le jeu et le serieux.
%
%  ET LA MENAGERIE EN CARTON EST A MOITIE UNE FAUNE DE JAVA, COMME
%  DANS LA COLONNE INDONESIENNE : badak le rhinoceros (7), kethek le
%  singe (10), baya le crocodile (12), macan tutul le leopard (18),
%  lawa la chauve-souris (35), wit klapa les cocotiers de la bataille
%  navale (48). Mais le javanais y ajoute ce que l'indonesien n'a pas
%  : kethek et baya sont des mots javanais purs, et l'indonesien dit
%  monyet et buaya. Deux langues au meme endroit de la meme ile, et
%  elles ne nomment pas le meme singe du meme mot.
%
%  ET LES QUATRE POINTS CARDINAUX (c) SONT ENCORE PLUS JAVANAIS QU'EN
%  INDONESIEN : lor, kidul, wetan, kulon — et non utara, selatan,
%  timur, barat. Ce sont les mots qui nomment les rues, les quartiers
%  et les portes des kraton depuis mille ans : Malioboro est au nord,
%  et le sud de Yogyakarta est le kidul de la mer. La rose des vents
%  se dit « mata angin » dans les deux langues, l'oeil du vent ; mais
%  seul le javanais a garde ses propres points.
% ===================================================================

%%K t16-apar-1 apar
\VUsaut{4.0mm}
\VUtitre{15.0pt}{60}{GAMBAR No. 16}
\VUsaut{2.0mm}
\VUpk{13.2pt}{40}{Toko Gedhe. --- Pasar.}
\VUpk{13.2pt}{40}{Dolanan.}
\VUsaut{3.4mm}

%%K t16-01-1 p
1. --- Taun anyar wis cedhak. Toko-toko gedhe wis nyiyapake pajangan
\VUgras{dolanan} \textsuperscript{(1)} lan \VUgras{hadiah taun anyar}e.

%%K t16-01-2 p
Aku njaluk marang Mbah Kakung supaya ngeterake aku menyang pasar
kanggo ndeleng pameran kang endah iku, lan panjenengane kersa.

%%K t16-02-1 p
2. --- Kang kapisan kita mandheg ing ngarepe tameng pajangan gegaman
kang apik, kang isi \VUgras{bedhil}, \VUgras{topi kepi},
\VUgras{pedhang}, \VUgras{sabuk kopel} lan sapasang \VUgras{epolet,}
utawa \VUgras{helm}, \VUgras{klambi waja} \textsuperscript{(3)} lan
\VUgras{pistul} \textsuperscript{(4)}. Kabeh iku luwih narik kawigatenku
tinimbang \VUgras{lentera ajaib} \textsuperscript{(2)}.

%%K t16-tit-1 sub
\VUsaut{3.0mm}
\VUpk{10.2pt}{40}{Pemandangan kang endah.}
\VUsaut{2.6mm}

%%K t16-03-1 p
3. --- Ing bagean kang padha ana \VUgras{prajurit jaga}
\textsuperscript{(5)} ing pos jagane ; dheweke katon lagi njaga ing
ngarepe sakebon binatang saka karton. Sepira akehe kewan ing kono :
\VUgras{manuk nuri} \textsuperscript{(6)} ing pancikane, \VUgras{badak}
\textsuperscript{(7)} kang ana sungune ing irunge, \VUgras{rusa kutub}
\textsuperscript{(8)}, \VUgras{manuk unta} \textsuperscript{(9)},
\VUgras{kethek} \textsuperscript{(10)} kang lagi nglethus kacang,
\VUgras{rubah} \textsuperscript{(11)} kang nggawa lunga pitik,
\VUgras{baya} \textsuperscript{(12)} kang mangapake rahange kang gedhe
temen, \VUgras{unta} \textsuperscript{(13)}, \VUgras{asu greyhound}
\textsuperscript{(14)} kang luwes, \VUgras{macan kumbang}
\textsuperscript{(15)}, banjur kewan loro kang ora dakkenal lan kang
jarene \VUgras{cerpelai} \textsuperscript{(16)} lan \VUgras{hiena}
\textsuperscript{(17)}. Ing kono uga ana \VUgras{macan tutul}
\textsuperscript{(18)} lan \VUgras{asu ajag} \textsuperscript{(21)} ; pas
ing cedhake ana \VUgras{tikus} \textsuperscript{(19)} cilik kang lucu lan
\VUgras{tikus cilik} \textsuperscript{(20)} karet kang cilik banget.
Pungkasane \VUgras{kuda nil} \textsuperscript{(22)} lan \VUgras{unta
punuk siji} \textsuperscript{(23)} kang bongkok nutup klumpukan iku. Aku
mesthi bakal kari ing kono nganti pirang-pirang jam yen ing sandhinge
ora ana perkemahan kang endah kang kebak \VUgras{serdadu timah}
\textsuperscript{(29)} kang narik kawigatenku.

%%K t16-tit-2 sub
\VUsaut{4.0mm}
\VUpk{10.2pt}{40}{Prajurit lan perang.}
\VUsaut{2.8mm}

%%K t16-04-1 p
4. --- Mbah Kakung, kang dadi \VUgras{veteran panyandhang disabilitas}
\textsuperscript{(24)} lan kang kudu ninggal tentara merga
\VUgras{tatu} kang njalari panjenengane nampa \VUgras{medali militer},
remen banget nyariyosaken marang aku \VUgras{kampanye perang} kang
tau dilakoni. Panjenengane bungah nerangake maneka gerakane pasukan
mungil iku. Sagerombolan prajurit tenanan kang lagi nyowani pasar —
\VUgras{prajurit zeni} \textsuperscript{(25)}, \VUgras{prajurit pemburu
Alpen} \textsuperscript{(26)} kang sirahe nganggo \VUgras{baret},
\VUgras{sersan mayor} \textsuperscript{(27)} infanteri lan \VUgras{sersan
artileri} \textsuperscript{(28)} kang \VUgras{pita pangkat} mas ing lengen
klambine — mandheg ing cedhak kita kanggo ngrungokake katrangan iku.

%%K t16-05-1 p
5. --- Mbah Kakung, kanthi bangga banget merga kagungan pamireng kang
mencorong temen, sanalika nyusun rancangan perang kang jangkep.

%%K t16-05-2 p
Kutha kang dibentengi, karo \VUgras{tembok kutha} lan \VUgras{parit}e,
dikepung \VUgras{tentara mungsuh}, kang, sawise \VUgras{pemboman} kang
suwe, siyap \VUgras{nyerbu}. Nanging \VUgras{sing kekepung}, merga
kedhesek luwe, nyoba \VUgras{metu nyerang} menyang \VUgras{perkemahan}
para \VUgras{pengepung}. Sawise mundurake \VUgras{serangan} iku ing
\VUgras{peperangan} kang sengit ing sakubenge \VUgras{tarub},
para pengepung kang \VUgras{menang} nglepasake \VUgras{skuadron}e
kanggo mburu para panyerang kang \VUgras{kalah}. Wong-wong iku
\VUgras{mundur} karo nutupi \VUgras{medan perang} nganggo \VUgras{kang
padha tiwas} lan \VUgras{kang padha tatu}. \VUgras{Para tukang usung
tandhu} nggawa kang keri iku menyang \VUgras{ambulans} supaya dirawat
dening \VUgras{ahli bedah}.

%%K t16-05-3 p
Sanajan ana panglawan kang kendel saka sawetara \VUgras{batalion} kang
wis mbentuk \VUgras{formasi pesagi}, \VUgras{kekalahan} iku sampurna ;
kekarene tentara \VUgras{mlayu} karo ninggal akeh \VUgras{tawanan}.
Mungsuh arep nyampurnakake \VUgras{kemenangan}e kanthi nduweni kutha
iku. Bisa uga iku bakal dadi pungkasane kampanye, lan, sawise
\VUgras{gencatan senjata}, wong bakal bisa nandhatangani
\VUgras{perjanjian damai}.

%%K t16-tit-3 sub
\VUsaut{4.4mm}
\VUpk{10.2pt}{40}{Hadiah kanggo taun anyar.}
\VUsaut{2.8mm}

%%K t16-06-1 p
6. --- Para prajurit enom, kang ngguyu ngrungokake critane sang veteran
kang endah, njaluk katrangan liyane marang panjenengane. Dene aku, aku
ngubengi gerai kanggo nyawang \VUgras{gandhewa silang}
\textsuperscript{(32)} kang endah lan \VUgras{gandhewa}
\textsuperscript{(33)} karo \VUgras{panah}e \textsuperscript{(31)}. Ing
kono aku nemu klumpukan kewan liyane : \VUgras{manuk kicau}
\textsuperscript{(34)} loro, yaiku \VUgras{manuk bulbul} lan
\VUgras{manuk cikrak} otomatis kang ngoceh sakayange, lan sarangkeyan
kewan aneh : \VUgras{lawa} \textsuperscript{(35)}, \VUgras{ula boa}
\textsuperscript{(36)}, \VUgras{kura-kura} \textsuperscript{(37)},
\VUgras{asu segara} \textsuperscript{(38)}, \VUgras{iwak paus}
\textsuperscript{(39)} lan \VUgras{iwak hiu} \textsuperscript{(40)} saka
karet tiup.

%%K t16-07-1 p
7. --- Aku ora suwe mandheg nyawang kewan iku, amarga aku weruh apa
kang tak-impekake : \VUgras{panggung boneka} \textsuperscript{(41)} kang
endah banget karo \VUgras{dekorasi} kang mewah lan \VUgras{tirai} abang
kang nengsemake. Ing \VUgras{panggung}, \VUgras{aktor} loro katon lagi
main \VUgras{peran} ing \VUgras{lakon} kang narik banget, amarga para
\VUgras{penonton} cilik saka karton kang dipasang ing \VUgras{balkon}
utawa kang lungguh ing \VUgras{kursi utama} lan ing \VUgras{lantai
ngisor} ing mburine \VUgras{dirigen}, padha keplok-keplok kanthi
bungah.

%%K t16-07-2 p
Emane, panggung iku regane larang banget.

%%K t16-08-1 p
8. --- Kejaba iku, milih dolanan iku angel, amarga ing etalase numpuk
sakabehing dolanan : \VUgras{Harlekin} \textsuperscript{(42)},
\VUgras{Pulcinella} \textsuperscript{(43)}, \VUgras{Pierrot}
\textsuperscript{(44)}, \VUgras{manuk kuwuk} \textsuperscript{(45)} kang
ngebutake elare, \VUgras{manuk ireng} \textsuperscript{(46)} kang nyuling,
\VUgras{asu pudel} kang nggonggong, \VUgras{fonograf}
\textsuperscript{(47)} lan \VUgras{dolanan katrampilan.} Salah sijine
dolanan iki nyenengake banget aku : dolanan iku

%%K t16-noto-1 noto
\VUnotes{143.2mm}{%
\textit{(*) Delengen gambar No. 6.}%
}

%%K t16-08-1 p suite
nggambarake \VUgras{peperangan laut} \textsuperscript{(48)} panggonan
\VUgras{kapal} diserang \VUgras{bajak laut} ing cedhake negara kang
ditanduri \VUgras{wit klapa.}

%%K t16-09-1 p
9. --- Aku bisa nyawang kabeh kaelokan iku kanthi gampang banget,
amarga mung sethithik wong ing toko, meh mung sawetara wong kang
ndelok-ndelok, \VUgras{prajurit kavaleri} \textsuperscript{(49)} siji lan
\VUgras{juru tagih} \textsuperscript{(50)} kang masrahake \VUgras{wesel}
ing \VUgras{kasir utama} \textsuperscript{(51)}.

%%K t16-10-1 p
10. --- Aku takon marang Mbah Kakung bab gunane buku-buku kang
diselehake ing rak ing mburine \VUgras{kasir wadon} iku ; panjenengane
ngendika yen buku-buku iku \VUgras{buku pembukuan}.

%%K t16-tit-4 sub
\VUsaut{3.6mm}
\VUpk{10.2pt}{40}{Ndelok-ndelok keliling toko.}
\VUsaut{2.8mm}

%%K t16-11-1 p
11. --- Nalika ninggal bagean dolanan, kita menyang bagean pelana
kanggo nyawang sedhela \VUgras{pelana} \textsuperscript{(53)},
\VUgras{sanggurdi} \textsuperscript{(54)}, \VUgras{tali kendhali}
\textsuperscript{(55)}, \VUgras{wesi kendhali} \textsuperscript{(56)} lan
\VUgras{cemeti}. Nalika \VUgras{kepala bagean} \textsuperscript{(57)}
menehi sawetara katrangan marang Mbah Kakung, aku lunga nyawang
pajangan \VUgras{porselen} lan \VUgras{kristal} \textsuperscript{(58)},
panggonan ana \VUgras{mangkok salad} kang apik lan \VUgras{teko}
\textsuperscript{(59)} kang alus.

%%K t16-12-1 p
12. --- Kanggo munggah menyang lantai kapisan, kita liwat ing mburine
etalase \VUgras{korset} \textsuperscript{(60)}, ing sandhinge bagean kaos
kaki, panggonan dipajang \VUgras{kathok jero} \textsuperscript{(63)} kang
apik banget. Ing cedhake, \VUgras{pramuniaga wadon}
\textsuperscript{(61)} nggawe \VUgras{wong tuku wadon}
\textsuperscript{(62)} milih \VUgras{kain} \textsuperscript{(64)} sutra
kang endah.

%%K t16-12-2 p
Karo munggah andha, aku weruh yen toko iku direngga \VUgras{lampion}
\textsuperscript{(65)} Jepang, \VUgras{payung} \textsuperscript{(66)} lan
\VUgras{wit palem} \textsuperscript{(70)} kang gedhe banget.

%%K t16-13-1 p
13. --- Ing lantai kapisan ora akeh kang bisa didelok ; mung perabot
(*), \VUgras{brankas} \textsuperscript{(67)}, \VUgras{lemari laci}
\textsuperscript{(68)} lan \VUgras{kreta bayi} \textsuperscript{(69)}.
Nanging pemandangan toko sakabehe \VUgras{nyenengake} banget.

%%K t16-14-1 p
14. --- Ing sangisore sikil kita, aku weruh alat-alat musik :
\VUgras{harpa} \textsuperscript{(71)}, \VUgras{gitar}
\textsuperscript{(72)}, \VUgras{mandolin} \textsuperscript{(73)},
\VUgras{kornet piston} \textsuperscript{(74)}, \VUgras{klarinet}
\textsuperscript{(75)}, biola karo \VUgras{penggesek}e
\textsuperscript{(76)}, lan malah \VUgras{genderang gedhe}
\textsuperscript{(77)}. Rada adoh sethithik, ing bagean alat tulis,
\VUgras{mahasiswa} \textsuperscript{(78)} lagi nawar map buku tulis marang
\VUgras{pramuniaga} \textsuperscript{(79)}. Ing cedhake wong loro iku, aku
weruh \VUgras{buku abjad} \textsuperscript{(80)} kang apik.

%%K t16-14-2 p
Ing tengahing toko ngadeg \VUgras{wit Natal} \textsuperscript{(81)} kang
dhuwur, kebak lilin, pernik lan sakabehing dolanan : \VUgras{jaran
kayu} \textsuperscript{(82)}, \VUgras{boneka} \textsuperscript{(83)}, lan
sapanunggalane.

%%K t16-14-3 p
\VUgras{Bagean parfum} \textsuperscript{(84)} karo \VUgras{banyu
pembersih untu}ne lan maneka \VUgras{lenga wangi} nyebar ambu kang
wangi banget, kang munggah tekan kita.

%%K t16-tit-5 sub
\VUsaut{3.4mm}
\VUpk{10.2pt}{40}{Kita tuku sawetara.}
\VUsaut{2.8mm}

%%K t16-15-1 p
15. --- Sarehne Mbah Kakung wiwit ngraos wektune kesuwen, kita mudhun
liwat andha gedhe. Panjenengane mandheg sedhela ing ngarepe
\VUgras{papan astronomi} \textsuperscript{(85)} kang nggambarake, kandhane
marang aku, \VUgras{komet} \textsuperscript{(a)}, \VUgras{grahana rembulan
lan srengenge} \textsuperscript{(b)}, mata angin karo \VUgras{papat
kiblat baku} \textsuperscript{(c)} \VUgras{(Lor, Kidul, Wetan lan
Kulon)}, peta \VUgras{belahan bumi} \textsuperscript{(d)} loro,
\VUgras{planet} \textsuperscript{(e)} lan \VUgras{tata surya}
\textsuperscript{(f)}. Aku ora mudheng apa-apa saka kabeh iku, lan aku
luwih kepencut marang seragam berpita \VUgras{bocah pengantar barang}
\textsuperscript{(86)} kang liwat, lan marang \VUgras{kipas}
\textsuperscript{(87)} endah kang ana ing cedhakku, sandhinge
\VUgras{dompet dhuwit receh} \textsuperscript{(88)} lan \VUgras{dompet}
\textsuperscript{(89)}.

%%K t16-16-1 p
16. --- Aku uga gumun nyawang ing gerai \VUgras{pena}
\textsuperscript{(90)} lan \VUgras{gesper} \textsuperscript{(91)} sabuk,
sarta \VUgras{kembang tiruan} \textsuperscript{(94)} kang endah kang
ditata luwes ing kranjang. \VUgras{Kembang violet}, \VUgras{kembang
stok}, \VUgras{kembang hiasin} \textsuperscript{(95)}, \VUgras{kembang
geranium} \textsuperscript{(96)}, \VUgras{kembang bakung}
\textsuperscript{(97)} lan \VUgras{kembang nasturtium}
\textsuperscript{(98)} ditiru kanthi becik banget.

%%K t16-tit-6 sub
\VUsaut{4.4mm}
\VUpk{10.2pt}{40}{Hadiah saka Mbah Kakung.}
\VUsaut{2.8mm}

%%K t16-17-1 p
17. --- Aku nyuwun marang Mbah Kakung supaya nukokake aku salah sijine
\VUgras{sikat untu} \textsuperscript{(92)} kang ana ing kothak cedhak
\VUgras{jepit rambut} \textsuperscript{(93)}. Panjenengane kersa, lan aku
dhewe lunga mbayar blanjanku ing kasir, panggonan wong lagi nyiyapake
\VUgras{kiriman} \textsuperscript{(100)} kang gedhe kanggo
\VUgras{pelanggan} \textsuperscript{(99)}.

%%K t16-18-1 p
18. --- Dumadakan ana rame cilik ing antarane wong-wong kang mandheg
cedhak \VUgras{barang-barang prunggu} \textsuperscript{(101)} seni.
\VUgras{Maling} \textsuperscript{(102)} lagi wae kecekel nalika lagi
nyolong dening \VUgras{pengawas} \textsuperscript{(103)}, nalika dheweke
nyoba nguras duweke wong. Wong enggal ngundang polisi kanggo
\VUgras{nyekel} dheweke, lan pas nalika kita metu, si pelaku lagi
digiring menyang pakunjaran.

\VUsaut{6.0mm}
\VUfilet{22mm}
